\documentclass[11pt]{article}

\usepackage[margin=1in]{geometry}
\usepackage{amsmath,amssymb,amsthm,mathtools}
\usepackage{enumitem}
\usepackage{xcolor}
\usepackage[colorlinks=true,linkcolor=blue,citecolor=blue,urlcolor=blue]{hyperref}

\newtheorem{theorem}{Theorem}
\newtheorem{lemma}[theorem]{Lemma}
\newtheorem{corollary}[theorem]{Corollary}
\theoremstyle{definition}
\newtheorem{definition}[theorem]{Definition}
\newtheorem{remark}[theorem]{Remark}

\newcommand{\TW}{Thatcher--Wright}
\newcommand{\Nat}{\mathbb N}
\newcommand{\powfin}{\mathcal P_{\mathrm{fin}}}
\newcommand{\SigmaSym}{\Sigma}
\newcommand{\rank}{\sigma}
\newcommand{\Term}[1]{T_{#1}}
\newcommand{\Aut}{\mathcal A}
\newcommand{\NAut}{\mathcal R}
\newcommand{\bh}{\mathrm{bh}}
\newcommand{\MSO}{\mathrm{MSO}}
\newcommand{\child}{\mathrm{child}}
\newcommand{\dom}{\mathrm{dom}}
\newcommand{\arity}{\mathrm{arity}}
\newcommand{\TODO}[1]{\textcolor{red}{\textbf{TODO:} #1}}

\title{Working TeX Extraction of Thatcher--Wright}
\author{For the Courcelle formalization project}
\date{}

\begin{document}
\maketitle

\begin{abstract}
This note extracts the part of Thatcher and Wright~\cite{TW} that is useful for
formalising the tree-automata ingredient in Courcelle's theorem.  The source PDF
is a scan with a noisy text layer, so the material below is a cleaned working
transcription/paraphrase rather than a line-by-line archival copy.  The goal is
to identify the definitions, lemmas, automaton constructions, and proof
obligations that should eventually become Lean objects.
\end{abstract}

\tableofcontents

\section{How This Fits Courcelle}

The proof of Courcelle's theorem needs the following finite-tree automata
principle.  This is the statement called Theorem~26 in
\texttt{lecture\_note\_expanded.tex}.

\begin{theorem}[Lecture-note Theorem 26: effective \(\MSO\)--tree-automaton
correspondence]\label{thm:mso-tree-automata}\label{thm:lecture-note-26}
Let \(\Sigma\) be a finite alphabet.  From an \(\MSO\) sentence \(\Phi\) over
\[
  \tau_\Sigma=\{\child_1,\child_2\}\cup\{P_a:a\in\Sigma\}
\]
one can compute a deterministic bottom-up finite tree automaton
\(\Aut_\Phi\) such that, for every \(\Sigma\)-tree, \(\Aut_\Phi\) accepts it iff
it satisfies \(\Phi\).  Running \(\Aut_\Phi\) takes time linear in the number of
nodes.
\end{theorem}

\TW{} prove the needed ingredients in a slightly different language.
Their ``generalized finite automata'' are finite algebras over a ranked
signature; their inputs are finite terms.  This is exactly the bottom-up
ranked-tree automata view.  Sections~2--3 give finite automata theory for terms:
determinisation, Boolean closure, projection, inverse projection, replacement,
emptiness, and regularity.  Section~4 performs the logic-to-automata induction
for weak monadic second-order logic over the binary successor tree.

For our Courcelle notes, the encoding of finite subsets of \(N_2\) in
\S\,4 can mostly be replaced by labelled input trees directly.  The core proof
pattern is still the same: atomic formulas are recognizable, and recognizable
relations are closed under conjunction, negation, and existential set
quantification.

To connect the displayed theorem with \TW{}'s ranked-term formalism, encode a
\(\Sigma\)-tree as a full binary ranked term by adding a fresh nullary symbol
\(\bot\).  Each real node labelled \(a\in\Sigma\) becomes a binary symbol
\(a(t_1,t_2)\); an absent child is represented by \(\bot\).  The original
relations \(\child_1,\child_2\) and predicates \(P_a\) are first-order definable
on the non-\(\bot\) part of this padded term.  The \(\MSO\)-to-recognizable
induction below therefore produces a bottom-up automaton for the padded ranked
terms.  Determinization gives the deterministic automaton in
Theorem~\ref{thm:mso-tree-automata}, and its run is linear because it performs
one transition per node of a padded term whose size is \(O(|T|)\).

\section{Generalized Automata}

\begin{definition}[Species and terms, TW \S2]
A \emph{species} is a pair \((\SigmaSym,\rank)\), where \(\SigmaSym\) is a set of
function symbols and \(\rank:\SigmaSym\to\Nat\) assigns each symbol its arity.
Write
\[
  \SigmaSym_n=\rank^{-1}(\{n\}).
\]
The set \(\Term{\SigmaSym}\) of finite terms is the least set such that
\(\SigmaSym_0\subseteq\Term{\SigmaSym}\) and, if \(f\in\SigmaSym_n\) and
\(t_1,\ldots,t_n\in\Term{\SigmaSym}\), then \(f(t_1,\ldots,t_n)\in
\Term{\SigmaSym}\).
\end{definition}

\begin{definition}[Deterministic automata]
A deterministic \(\SigmaSym\)-automaton is a finite
\(\SigmaSym\)-algebra
\[
  \Aut=(A,\alpha),
\]
where \(A\) is a finite set of states and, for each \(f\in\SigmaSym_n\),
\(\alpha_f:A^n\to A\).  The constants \(f\in\SigmaSym_0\) are initial states.
The bottom-up run is the unique homomorphism
\[
  h_\Aut:\Term{\SigmaSym}\to A
\]
defined by
\[
  h_\Aut(f(t_1,\ldots,t_n))
  =\alpha_f(h_\Aut(t_1),\ldots,h_\Aut(t_n)).
\]
For final states \(A_F\subseteq A\), the behaviour is
\[
  \bh_\Aut(A_F)=\{t\in\Term{\SigmaSym}:h_\Aut(t)\in A_F\}.
\]
\end{definition}

\begin{definition}[Nondeterministic automata]
A nondeterministic \(\SigmaSym\)-automaton is a finite relational algebra
\[
  \NAut=(R,\rho),
\]
where, for \(f\in\SigmaSym_n\), \(\rho_f\subseteq R^n\times R\).  Its run returns
a set of possible states:
\[
  h_\NAut(f(t_1,\ldots,t_n))
  =
  \{r:\exists r_1,\ldots,r_n,\ r_i\in h_\NAut(t_i)
      \text{ and } \rho_f(r_1,\ldots,r_n,r)\}.
\]
The behaviour with final states \(R_F\subseteq R\) is
\[
  \bh_\NAut(R_F)=\{t:h_\NAut(t)\cap R_F\neq\emptyset\}.
\]
\end{definition}

\begin{definition}[Recognizable sets]
A set \(U\subseteq\Term{\SigmaSym}\) is \emph{recognizable} if
\(U=\bh_\Aut(A_F)\) for some finite deterministic or nondeterministic
\(\SigmaSym\)-automaton.
\end{definition}

\begin{theorem}[TW Theorem 1: subset construction]\label{thm:tw-subset}
A set of terms is recognizable by a deterministic automaton iff it is
recognizable by a nondeterministic automaton.
\end{theorem}

\begin{proof}[Construction]
Every deterministic automaton is nondeterministic by replacing each transition
function by its graph.  More explicitly, if \(\Aut=(A,\alpha)\) is deterministic,
define a nondeterministic automaton on the same state set by
\[
  \rho_f(a_1,\ldots,a_n,a)
  \quad\Longleftrightarrow\quad
  \alpha_f(a_1,\ldots,a_n)=a
\]
for every \(f\in\SigmaSym_n\).  For \(n=0\) this says that the only possible
state at the constant \(f\) is \(\alpha_f\).  A structural induction on \(t\)
shows that the nondeterministic run is always the singleton
\[
  h_\NAut(t)=\{h_\Aut(t)\},
\]
so the same final states recognize the same language.

Conversely, let \(\NAut=(R,\rho)\) be a nondeterministic automaton.  Define the
deterministic powerset automaton
\[
  \widehat{\NAut}=(\mathcal P(R),\widehat\rho)
\]
as follows.  For \(f\in\SigmaSym_n\), set
\[
  \widehat\rho_f(S_1,\ldots,S_n)
  =
  \{r\in R:\exists r_1\in S_1,\ldots,\exists r_n\in S_n,\,
      \rho_f(r_1,\ldots,r_n,r)\}.
\]
When \(n=0\), the right hand side is read using the unique empty tuple:
\[
  \widehat\rho_f
  =
  \{r\in R:\rho_f(r)\}.
\]
Since \(R\) is finite, \(\mathcal P(R)\) is finite, so this is again a finite
automaton.  If \(R_F\) is the nondeterministic accepting set, take the accepting
states of the powerset automaton to be
\[
  \widehat R_F=\{S\subseteq R:S\cap R_F\neq\emptyset\}.
\]

It remains to prove that this construction has the same run semantics as the
original nondeterministic automaton.  We prove, by induction on terms, that
\[
  h_{\widehat{\NAut}}(t)=h_\NAut(t)       \tag{\(*\)}
\]
for every \(t\in\Term{\SigmaSym}\).

For the base case, let \(t=f\) with \(f\in\SigmaSym_0\).  By the definition of
the deterministic run in \(\widehat{\NAut}\),
\[
  h_{\widehat{\NAut}}(f)=\widehat\rho_f
  =
  \{r\in R:\rho_f(r)\}.
\]
This is exactly the definition of the nondeterministic run \(h_\NAut(f)\) at the
constant \(f\), so \((*)\) holds.

For the inductive step, let \(t=f(t_1,\ldots,t_n)\), and assume
\[
  h_{\widehat{\NAut}}(t_i)=h_\NAut(t_i)
  \qquad (1\le i\le n).
\]
Then
\begin{align*}
  h_{\widehat{\NAut}}(t)
  &=
  \widehat\rho_f\bigl(h_{\widehat{\NAut}}(t_1),\ldots,
    h_{\widehat{\NAut}}(t_n)\bigr) \\
  &=
  \{r\in R:\exists r_1\in h_{\widehat{\NAut}}(t_1),\ldots,
      \exists r_n\in h_{\widehat{\NAut}}(t_n),\,
      \rho_f(r_1,\ldots,r_n,r)\} \\
  &=
  \{r\in R:\exists r_1\in h_\NAut(t_1),\ldots,
      \exists r_n\in h_\NAut(t_n),\,
      \rho_f(r_1,\ldots,r_n,r)\} \\
  &=
  h_\NAut(f(t_1,\ldots,t_n)).
\end{align*}
The third equality is the induction hypothesis, and the last equality is the
definition of \(h_\NAut\).  This proves \((*)\) for all terms.

Finally, for every term \(t\),
\[
  t\in\bh_{\widehat{\NAut}}(\widehat R_F)
  \Longleftrightarrow
  h_{\widehat{\NAut}}(t)\cap R_F\neq\emptyset
  \Longleftrightarrow
  h_\NAut(t)\cap R_F\neq\emptyset
  \Longleftrightarrow
  t\in\bh_\NAut(R_F).
\]
Thus the powerset automaton recognizes exactly the same set of terms as
\(\NAut\), completing the nontrivial direction.
\end{proof}

\begin{theorem}[TW Theorem 2: Boolean closure]\label{thm:tw-boolean}
If \(U,V\subseteq\Term{\SigmaSym}\) are recognizable, then so are
\(U\cap V\), \(U\cup V\), and \(\Term{\SigmaSym}\setminus U\).
\end{theorem}

\begin{proof}[Construction]
By Theorem~\ref{thm:tw-subset}, we may assume the recognizing automata are
deterministic.  Let
\[
  \Aut=(A,\alpha),\qquad \mathcal B=(B,\beta)
\]
be deterministic \(\SigmaSym\)-automata with final-state sets
\(A_F\subseteq A\) and \(B_F\subseteq B\), and suppose
\[
  \bh_\Aut(A_F)=U,\qquad \bh_{\mathcal B}(B_F)=V.
\]

\emph{Complement.}  Use the same automaton \(\Aut\), but replace the accepting
set \(A_F\) by \(A\setminus A_F\).  For every term \(t\),
\[
  t\in\bh_\Aut(A\setminus A_F)
  \Longleftrightarrow
  h_\Aut(t)\in A\setminus A_F
  \Longleftrightarrow
  h_\Aut(t)\notin A_F
  \Longleftrightarrow
  t\notin\bh_\Aut(A_F).
\]
Thus
\[
  \bh_\Aut(A\setminus A_F)=\Term{\SigmaSym}\setminus U.
\]
This is why determinism is useful here: each term has one state, so changing the
final set really complements the language.  For a nondeterministic automaton,
one must first determinize.

\emph{Intersection.}  Define the product automaton
\[
  \mathcal C=(A\times B,\gamma).
\]
For a constant \(f\in\SigmaSym_0\), put
\[
  \gamma_f=(\alpha_f,\beta_f)\in A\times B.
\]
For \(f\in\SigmaSym_n\) with \(n>0\), define
\[
  \gamma_f\bigl((a_1,b_1),\ldots,(a_n,b_n)\bigr)
  =
  \bigl(\alpha_f(a_1,\ldots,a_n),\beta_f(b_1,\ldots,b_n)\bigr).
\]
Equivalently, the same formula also covers \(n=0\) if one reads the empty tuple
correctly.

The crucial fact is that the product automaton computes both original runs at
the same time:
\[
  h_{\mathcal C}(t)=\bigl(h_\Aut(t),h_{\mathcal B}(t)\bigr)       \tag{\(*\)}
\]
for every \(t\in\Term{\SigmaSym}\).  We prove this by structural induction on
\(t\).

If \(t=f\) is a constant, then
\[
  h_{\mathcal C}(f)=\gamma_f=(\alpha_f,\beta_f)
  =
  \bigl(h_\Aut(f),h_{\mathcal B}(f)\bigr).
\]
For the inductive step, let \(t=f(t_1,\ldots,t_n)\), and assume \((*)\) holds
for each \(t_i\).  Then
\begin{align*}
  h_{\mathcal C}(t)
  &=
  \gamma_f\bigl(h_{\mathcal C}(t_1),\ldots,h_{\mathcal C}(t_n)\bigr) \\
  &=
  \gamma_f\bigl((h_\Aut(t_1),h_{\mathcal B}(t_1)),\ldots,
    (h_\Aut(t_n),h_{\mathcal B}(t_n))\bigr) \\
  &=
  \bigl(\alpha_f(h_\Aut(t_1),\ldots,h_\Aut(t_n)),
        \beta_f(h_{\mathcal B}(t_1),\ldots,h_{\mathcal B}(t_n))\bigr) \\
  &=
  \bigl(h_\Aut(f(t_1,\ldots,t_n)),
        h_{\mathcal B}(f(t_1,\ldots,t_n))\bigr).
\end{align*}
This proves \((*)\).

Now take final states \(A_F\times B_F\) in the product automaton.  Using
\((*)\), for every term \(t\),
\begin{align*}
  t\in\bh_{\mathcal C}(A_F\times B_F)
  &\Longleftrightarrow
  h_{\mathcal C}(t)\in A_F\times B_F \\
  &\Longleftrightarrow
  h_\Aut(t)\in A_F \text{ and } h_{\mathcal B}(t)\in B_F \\
  &\Longleftrightarrow
  t\in U \text{ and } t\in V.
\end{align*}
Therefore
\[
  \bh_{\mathcal C}(A_F\times B_F)=U\cap V.
\]
This proves that the product automaton with accepting set \(A_F\times B_F\) is
the desired automaton for intersection.

\emph{Union.}  Use the same product automaton \(\mathcal C\), but take accepting
states
\[
  (A_F\times B)\cup(A\times B_F).
\]
Again using \((*)\), for every term \(t\),
\begin{align*}
  t\in\bh_{\mathcal C}\bigl((A_F\times B)\cup(A\times B_F)\bigr)
  &\Longleftrightarrow
  h_\Aut(t)\in A_F \text{ or } h_{\mathcal B}(t)\in B_F \\
  &\Longleftrightarrow
  t\in U \text{ or } t\in V.
\end{align*}
Hence this product automaton recognizes \(U\cup V\).  This also follows from
complement and intersection by De Morgan's laws, but the direct final-state set
makes the automaton construction explicit.
\end{proof}

\section{Projection}

\begin{definition}[Projection of species]
Let \((\SigmaSym,\rank)\) and \((\Gamma,\rank')\) be species.  A map
\(\pi:\SigmaSym\to\Gamma\) is a projection if it preserves arities:
\[
  \rank'(\pi(f))=\rank(f).
\]
It extends to terms by
\[
  \pi(f(t_1,\ldots,t_n))
  =
  \pi(f)(\pi(t_1),\ldots,\pi(t_n)).
\]
\end{definition}

\begin{theorem}[TW Theorem 3: closure under projection]\label{thm:tw-projection}
If \(U\subseteq\Term{\SigmaSym}\) is recognizable and
\(\pi:\Term{\SigmaSym}\to\Term{\Gamma}\) is induced by an arity-preserving symbol
map, then \(\pi(U)\subseteq\Term{\Gamma}\) is recognizable.
\end{theorem}

\begin{proof}[Construction]
Let \(\Aut=(A,\alpha)\) recognize \(U\).  Build a nondeterministic
\(\Gamma\)-automaton on the same state set.  For a \(\Gamma\)-symbol \(g\), allow
all transitions of \(\Aut\) labelled by symbols \(f\) with \(\pi(f)=g\).  For a
constant \(g\), the initial state set is the set of \(\alpha_f\) over all
constants \(f\) with \(\pi(f)=g\).  The proof is by induction on terms: every
source run projects to a target run, and every target nondeterministic run can
be lifted to some source term.
\end{proof}

\begin{theorem}[TW Theorem 4: closure under inverse projection]
If \(U\subseteq\Term{\Gamma}\) is recognizable and
\(\pi:\Term{\SigmaSym}\to\Term{\Gamma}\) is induced by an arity-preserving symbol
map, then \(\pi^{-1}(U)\subseteq\Term{\SigmaSym}\) is recognizable.
\end{theorem}

\begin{proof}[Construction]
Given a \(\Gamma\)-automaton \((A,\alpha)\), define a \(\SigmaSym\)-automaton by
using the transition \(\alpha_{\pi(f)}\) for each \(f\in\SigmaSym\).  An induction
on terms proves \(h(t)=h(\pi(t))\).
\end{proof}

\section{Replacement and Emptiness}

\begin{definition}[Subterm and depth]
The depth of a term is
\[
  d(c)=1 \quad(c\in\SigmaSym_0), \qquad
  d(f(t_1,\ldots,t_n))=1+\max_i d(t_i).
\]
Write \(t\succeq s\) when \(s\) occurs as a subterm of \(t\).
\end{definition}

\begin{lemma}[TW Lemma 5: replacement]\label{lem:tw-replacement}
Let \(\Aut\) be a deterministic \(\SigmaSym\)-automaton.  Suppose \(t'\) is
obtained from \(t\) by replacing an occurrence of \(t_2\) by \(t_1\).  If
\[
  h_\Aut(t_1)=h_\Aut(t_2),
\]
then
\[
  h_\Aut(t')=h_\Aut(t).
\]
\end{lemma}

\begin{proof}
Induct on the context around the replaced subterm.  At the replaced node the two
states are equal by hypothesis; every ancestor computes the same transition from
the same tuple of child states.
\end{proof}

\begin{corollary}[TW corollary to Lemma 5]
The same conclusion holds when finitely many pairwise incomparable occurrences
are replaced by terms with the same automaton states.
\end{corollary}

\begin{lemma}[TW Lemma 6: bounded-depth witness]
Let \(\Aut\) have \(n\) states.  If some term \(t\) has \(h_\Aut(t)=a\), then
there is a term \(t'\) with \(d(t')\le n\) and \(h_\Aut(t')=a\).
\end{lemma}

\begin{proof}
If \(d(t)>n\), choose a root-to-leaf chain of more than \(n\) subterms.  Two
subterms on the chain have the same automaton state.  Replacing the upper one by
the lower one strictly reduces the term while preserving the root state by
Lemma~\ref{lem:tw-replacement}.  Iterate until depth is at most \(n\).
\end{proof}

\begin{theorem}[TW Theorem 7: decidable emptiness]
Given a finite \(\SigmaSym\)-automaton \(\Aut\) and final states \(A_F\), it is
decidable whether \(\bh_\Aut(A_F)=\emptyset\).
\end{theorem}

\begin{proof}
Enumerate the finitely many terms of depth at most the number of states and test
whether any reaches a final state.  The bounded-depth lemma proves completeness
of this finite search.
\end{proof}

\section{Regularity}

\begin{definition}[Complex product and closure, TW \S3]
Fix an individual symbol \(\lambda\in\SigmaSym_0\).  For
\(U,V\subseteq\Term{\SigmaSym}\), the complex product
\[
  U\cdot_\lambda V
\]
is the set of all terms obtained from a term \(u\in U\) by replacing every
occurrence of \(\lambda\) by an independently chosen term of \(V\).

The \(\lambda\)-closure \(U^{*_\lambda}\) is defined by
\[
  X_0=\{\lambda\}, \qquad
  X_{n+1}=X_n\cup(U\cdot_\lambda X_n), \qquad
  U^{*_\lambda}=\bigcup_{n\ge0}X_n.
\]
\end{definition}

\begin{definition}[Regular sets of terms]
A set \(U\subseteq\Term{\SigmaSym}\) is regular if it belongs to the least class
containing all finite sets and closed under finite union, complex products, and
closures, after allowing the signature to be extended by additional constant
symbols but not by additional positive-arity symbols.
\end{definition}

\begin{theorem}[TW Theorem 8: recognizable iff regular]
For every finite species, \(U\subseteq\Term{\SigmaSym}\) is recognizable iff it
is regular.
\end{theorem}

\begin{proof}[Proof structure]
\TW{} split the proof into analysis and synthesis.

Analysis (TW Theorem~9): every recognizable set is regular.  The proof adds one
fresh constant \(\lambda_a\) for each automaton state \(a\).  Regular sets
\(T^k_i\) are built by induction so that \(T^k_i\) contains exactly the terms
whose partial automaton computation reaches state \(a_i\) using only the first
\(k\) states as non-initial transition targets.  The induction step is expressed
using complex product and closure.

Synthesis: finite sets are recognizable (Lemma~10), recognizable sets are closed
under union (Theorem~2), complex product (Lemma~11), and closure (Lemma~12).
\end{proof}

\begin{lemma}[TW Lemma 10]
Every finite subset of \(\Term{\SigmaSym}\) is recognizable.
\end{lemma}

\begin{proof}[Idea]
Construct an automaton recognizing a singleton term by matching its finite tree
shape bottom-up, then use Boolean closure for finite unions.
\end{proof}

\begin{lemma}[TW Lemma 11: closure under complex product]
If \(U,V\subseteq\Term{\SigmaSym}\) are recognizable, then
\[
  U\cdot_\lambda V
\]
is recognizable for every \(\lambda\in\SigmaSym_0\).
\end{lemma}

\begin{proof}[Construction idea]
Given automata for \(U\) and \(V\), construct a nondeterministic automaton that
either simulates the \(U\)-automaton or, when it reaches a \(\lambda\)-hole,
switches to a complete run of the \(V\)-automaton.  If the inserted \(V\)-term is
accepted, the surrounding run may continue as though the \(\lambda\)-constant had
been read.  The proof uses a simultaneous induction: one invariant describes
states arising from the \(U\)-side, the other says the \(V\)-side simulation is
faithful.
\end{proof}

\begin{lemma}[TW Lemma 12: closure under complex closure]
If \(U\subseteq\Term{\SigmaSym}\) is recognizable, then \(U^{*_\lambda}\) is
recognizable.
\end{lemma}

\begin{proof}[Construction idea]
Use a nondeterministic automaton that simulates the automaton for \(U\), and
whenever a completed subterm is accepted as an element of \(U\), it may collapse
that subterm to the state corresponding to the constant \(\lambda\).  This
implements iterated substitution at \(\lambda\).
\end{proof}

\begin{definition}[Truncation, TW \S3]
For \(U,V\subseteq\Term{\SigmaSym}\), define
\[
  U/_{\lambda}V
  =
  \{t: ( \{t\}\cdot_\lambda V )\cap U\neq\emptyset\}.
\]
Intuitively, \(t\) is obtained from some term of \(U\) by replacing selected,
non-overlapping subterms from \(V\) by \(\lambda\), with every occurrence of
\(\lambda\) covered by one of the removed subterms.
\end{definition}

\begin{lemma}[TW Lemma 13: closure under truncation]\label{lem:tw-truncation}
If \(U\subseteq\Term{\SigmaSym}\) is recognizable and
\(V\subseteq\Term{\SigmaSym}\) is arbitrary, then \(U/_{\lambda}V\) is
recognizable.
\end{lemma}

\begin{proof}[Construction]
Let \(\Aut=(A,\alpha)\) recognize \(U\), and let
\[
  A_V=\{h_\Aut(t):t\in V\}.
\]
Construct a nondeterministic automaton with the same ordinary transitions as
\(\Aut\), but with the constant \(\lambda\) allowed to start in any state of
\(A_V\).  The proof is by induction on terms and establishes:
\[
  t'\in t\cdot_\lambda V \Rightarrow h_\Aut(t')\in h(t),
\]
and conversely, every state in the new run of \(t\) is realized by some
\(t'\in t\cdot_\lambda V\).
\end{proof}

\section{Weak MSO Over the Binary Successor Tree}

\TW{} apply the automata theory to the weak monadic second-order theory of
multiple successors.  For Courcelle, this is the important logic-to-automata
part of the paper.

\begin{definition}[The tree \(N_2\)]
Let \(N_2=\{1,2\}^*\), with empty word \(\epsilon\), and successor functions
\[
  r_i(w)=wi \qquad (i=1,2).
\]
The weak MSO language has first-order variables ranging over \(N_2\), set
variables ranging over finite subsets of \(N_2\), equality, membership, and
binary successor predicates \(R_i(x,y)\), interpreted as \(r_i(x)=y\).
\end{definition}

\begin{theorem}[TW Theorem 14, Doner]
The weak monadic second-order theory of multiple successors is decidable.
\end{theorem}

\begin{proof}[TW proof plan]
It suffices to treat \(N_2\).  The proof has three steps:
\begin{enumerate}[leftmargin=*]
  \item translate away first-order variables by replacing them with singleton set
    variables;
  \item encode \(n\)-tuples of finite subsets of \(N_2\) as finite terms over a
    finite dyadic species;
  \item prove that every relation definable in the set-only language is
    recognizable under this encoding.
\end{enumerate}
Decidability then follows by effectively constructing the automaton associated
to a sentence and applying the emptiness theorem.
\end{proof}

\section{Eliminating First-Order Variables}

\begin{definition}[Set-only translation]
Replace each first-order variable \(x\) by a set variable \(X_x\), constrained to
be a singleton.  A convenient cleaned definition of singleton, using only
\(\subseteq\), is:
\[
  \mathrm{Empty}(A) \;:\!\!\Longleftrightarrow\; \forall B,\ A\subseteq B,
\]
\[
  \mathrm{Sing}(A) \;:\!\!\Longleftrightarrow\;
  \neg\mathrm{Empty}(A)\ \wedge\
  \forall B,\ B\subseteq A \to (\mathrm{Empty}(B)\vee A\subseteq B).
\]
The atomic translation is:
\[
  (x=y)^* \;=\; X_x\subseteq X_y\wedge\mathrm{Sing}(X_x)\wedge
  \mathrm{Sing}(X_y),
\]
\[
  (x\in A)^* \;=\; X_x\subseteq A\wedge\mathrm{Sing}(X_x),
\]
\[
  (R_i(x,y))^* \;=\; \widehat R_i(X_x,X_y),
\]
where \(\widehat R_i\) is the successor relation restricted to singleton sets.
The translation commutes with Boolean connectives and sends
\[
  (\exists x\,\varphi)^*
  =
  \exists X_x\,(\mathrm{Sing}(X_x)\wedge\varphi^*),
  \qquad
  (\exists A\,\varphi)^*=\exists A\,\varphi^*.
\]
\end{definition}

\begin{remark}
The PDF gives an equivalent formula for \(\mathrm{Sing}\), but the OCR in that
line is especially noisy.  The formula above is the same idea in a form better
suited to formalization.
\end{remark}

\section{Encoding Finite Subsets as Terms}

\begin{definition}[Single-track encoding]
Let \(\Sigma_1\) have one nullary symbol \(\lambda\) and two binary symbols
\(0,1\).  Define
\[
  c:\Term{\Sigma_1}\to \powfin(N_2)
\]
by
\[
  c(\lambda)=\emptyset,
\]
\[
  c(0(t_1,t_2))=1c(t_1)\cup 2c(t_2),
\]
\[
  c(1(t_1,t_2))=\{\epsilon\}\cup 1c(t_1)\cup 2c(t_2),
\]
where \(ic(t)=\{iw:w\in c(t)\}\).  Thus a node labelled \(1\) marks the
corresponding position as present in the finite subset of \(N_2\), and a node
labelled \(0\) leaves it absent.
\end{definition}

\begin{definition}[Canonical encoding]
The map \(c\) is onto but not one-to-one, since all-zero subtrees carry no set
information.  For a finite set \(A\subseteq N_2\), let \(e(A)\) be the smallest
term in \(c^{-1}(A)\), equivalently the unique representative with no subterm of
the form \(0(\lambda,\lambda)\).
\end{definition}

\begin{lemma}[TW Lemma 15]\label{lem:tw-encoding-one}
The set
\[
  U_1=\{e(A):A\in\powfin(N_2)\}
\]
of canonical encodings of finite subsets of \(N_2\) is recognizable.
\end{lemma}

\begin{proof}[TW proof idea]
\TW{} exhibit \(U_1\) as a regular set.  Let \(U'_1\) be the finite set of
one-step terms that can occur in a canonical encoding after a temporary hole
symbol is allowed.  Its complex closure generates all candidate terms, and
complex product with the empty set removes the temporary hole.  An induction on
term depth proves that the generated terms are exactly the canonical encodings.
Theorem~8 then gives recognizability.
\end{proof}

\begin{definition}[Multi-track encoding]
For \(n\) set variables, use the species \(\Sigma_n\) with one nullary symbol
\(\lambda\) and binary symbols \(\{0,1\}^n\).  The projection
\[
  \pi_i:\{0,1\}^n\to\{0,1\}
\]
selects the \(i\)-th track and extends to terms.  Define
\[
  c_n(t)=\bigl(c(\pi_1(t)),\ldots,c(\pi_n(t))\bigr).
\]
For \(\bar A=(A_1,\ldots,A_n)\), \(e(\bar A)\) is the smallest term in
\(c_n^{-1}(\bar A)\), equivalently the representative with no subterm
\[
  (0,\ldots,0)(\lambda,\lambda).
\]
\end{definition}

\begin{lemma}[TW Lemma 16]\label{lem:tw-encoding-n}
The set
\[
  U_n=\{e(A_1,\ldots,A_n):(A_i\in\powfin(N_2))\}
\]
of canonical encodings of \(n\)-tuples of finite subsets of \(N_2\) is
recognizable.
\end{lemma}

\begin{proof}[Idea]
Same proof as Lemma~\ref{lem:tw-encoding-one}, using the all-zero \(n\)-track
label in place of the single label \(0\).
\end{proof}

\section{Definable Relations Are Recognizable}

\begin{definition}[Recognizable relation]
A relation \(R\subseteq(\powfin(N_2))^n\) is recognizable if
\[
  e(R)=\{e(\bar A):\bar A\in R\}
  \subseteq\Term{\Sigma_n}
\]
is recognizable.
\end{definition}

\begin{theorem}[TW Theorem 17]\label{thm:tw-definable-recognizable}
Every relation definable in the set-only weak MSO language over \(N_2\) is
recognizable.
\end{theorem}

\begin{proof}[Structure]
Induct on formulas.  Lemma~\ref{lem:tw-atomic} handles atomic formulas.
Lemma~\ref{lem:tw-logical-closure} handles conjunction, negation, and existential
set quantification.  The construction is effective because each closure theorem
gives an explicit automaton construction.
\end{proof}

\begin{lemma}[TW Lemma 18: atomic formulas]\label{lem:tw-atomic}
The relations on \(\powfin(N_2)\) defined by atomic formulas in the set-only
language are recognizable.
\end{lemma}

\begin{proof}[Working formalization version]
The atomic relations are subset \(A\subseteq B\) and singleton-successor
\(\widehat R_i(A,B)\).

For \(A\subseteq B\), use the two-track alphabet.  A canonical term violates
subset exactly at a node whose label is \((1,0)\).  Thus the subset relation is
the recognizable set of canonical encodings that avoid label \((1,0)\); this can
be recognized by a two-state bottom-up automaton tracking whether a violation
has appeared.

For \(\widehat R_i(A,B)\), first require that both tracks are singletons and then
check that the unique \(B\)-position is the \(i\)-child of the unique
\(A\)-position.  \TW{} prove this through explicit regular expressions and
Theorem~8.  For Lean, a direct bottom-up automaton is probably cleaner: propagate
whether zero, one, or many marked nodes occur on each track, together with the
local parent/child relation needed for \(R_i\).
\end{proof}

\begin{lemma}[TW Lemma 19: logical closure]\label{lem:tw-logical-closure}
If \(R\) and \(S\) are recognizable relations, then the following are
recognizable:
\[
  R(\bar A,\bar B)\wedge S(\bar B,\bar C),\qquad
  \neg R(\bar A),\qquad
  \exists A_i\,R(A_1,\ldots,A_m,\bar B).
\]
\end{lemma}

\begin{proof}[Constructions]
For conjunction, pull both relations back to a common product alphabet by inverse
projection, then intersect.

For negation, complement the recognizable set inside all canonical encodings:
\[
  e(\neg R)=\bigl(\Term{\Sigma_n}\setminus e(R)\bigr)\cap U_n.
\]
Use Boolean closure and Lemma~\ref{lem:tw-encoding-n}.

For existential quantification, project away the quantified track.  The image may
contain noncanonical terms because deleting a track can introduce all-zero
subtrees.  \TW{} repair this by truncation:
\[
  e(\exists A\,R)
  =
  \bigl(\pi(e(R))/_{\lambda}\{(0,\ldots,0)(\lambda,\lambda)\}^{*_\lambda}\bigr)
  \cap U_{n-1}.
\]
The exact all-zero term depends on the remaining number of tracks.  The
important formal point is that projection plus truncation plus intersection with
canonical encodings implements existential set quantification.
\end{proof}

\begin{corollary}[Effective MSO-to-automata theorem over \(N_2\)]
For every weak MSO formula over the binary successor tree, one can effectively
construct a finite tree automaton accepting exactly the canonical encodings of
its satisfying finite-set assignments.
\end{corollary}

\begin{proof}
Translate away first-order variables, then apply
Theorem~\ref{thm:tw-definable-recognizable}.  Determinize with
Theorem~\ref{thm:tw-subset} if a deterministic automaton is desired.
\end{proof}

\section{Lean-Oriented Decomposition}

For the Courcelle formalization, this paper suggests the following modules or
proof blocks.
\begin{enumerate}[leftmargin=*]
  \item Ranked signatures and finite terms.
  \item Deterministic bottom-up tree automata and runs.
  \item Nondeterministic automata and subset construction.
  \item Boolean closure, projection, and inverse projection constructions.
  \item Atomic automata for labelled tree vocabularies.
  \item Formula induction compiling MSO formulas to automata.
  \item Optional: canonical-encoding machinery from \TW{} \S4 if we choose to
    formalize the \(N_2\)-subset version literally.
\end{enumerate}

The Courcelle proof itself can use a simpler finite-labelled-tree version of
\TW{}'s Section~4: assignments to MSO variables are represented as extra tracks
on the input tree, and existential quantification erases a track by projection.
In that setting, the truncation machinery is unnecessary unless we introduce a
separate canonicalization step.

\begin{thebibliography}{9}

\bibitem{TW}
James W. Thatcher and Jesse B. Wright.
\newblock Generalized finite automata theory with an application to a decision
problem of second-order logic.
\newblock \emph{Mathematical Systems Theory}, 2(1):57--81, 1968.

\end{thebibliography}

\end{document}
